\section{No-washout identification under arbitrary bounded adaptive baselines}
\label{sec:linear-time-protocol}
Put \(g_\beta(t)=\ell e^{tA_\beta}B=h_\beta'(t)\).
Product convergence gives uniform convergence of \(g_\beta\) on every
bounded time interval, and common stability gives the tail envelope
\(C_*e^{-\lambda_*t}\). Dominated convergence therefore proves
continuity of \(\beta\mapsto g_\beta\) into \(L^1(0,\infty)\), with
compact image. For all realized inputs with \(\|u\|_\infty\le U_*\),
\begin{equation}
 \sup_{t\ge0}|((g_\beta-g_\gamma)*u)(t)|
 \le U_*\|g_\beta-g_\gamma\|_1.
 \label{eq:policy-uniform-convolution}
\end{equation}
This holds for arbitrary switching and feedback histories, not just for a
particular controller.

Fix \(\tau_g>0\), \(T_0,U_0<\infty\), and \(\rho\in(0,1]\).
At stage \(i\), the controller first chooses a baseline segment of duration
\(r_i\in[\tau_g,\tau_g+T_0]\). A fresh exploration flag has probability
\(\rho\); on exploration draw a cell \(M_i\sim\omega\), otherwise the
controller selects any cell. The baseline during both the first segment and
the final probe is bounded by \(U_0\), and is fixed before the fresh sign.
Then add \(aS_i\) during the final interval \(t_{M_i}\), and draw readout
noise only at its end. Any baseline randomizers are included before the
sign. The parameter-independent policy may depend on the horizon. Thus
\begin{equation}
 \mathcal F_{i-1}\longrightarrow(r_i,u_i^0,E_i,M_i)
 \longrightarrow S_i\longrightarrow\xi_i\longrightarrow Y_i,
 \qquad \tau_g n\le\mathsf T_n\le\tau_+n,
 \quad\tau_+=\tau_g+T_0+T.
 \label{eq:linear-time-chronology}
\end{equation}
No washout or state reset is required. Write \(U_*=U_0+a\).

\subsection{A deterministic metric controlling every policy}
Define the compact metric
\begin{equation}
 d_*(\beta,\gamma)=U_*\|g_\beta-g_\gamma\|_1
                  +a\|h_\beta-h_\gamma\|_{C([0,T])}.
 \label{eq:policy-controlling-metric}
\end{equation}
Continuity follows from the compact response images. Injectivity follows
from boundary reconstruction, so this metric induces the product topology.
For \(f_i(\beta)=m_i^\beta-m_i^{\beta_0}\), where \(m_i\) is the mean
from zero initial state under the complete realized input, linearity gives
\begin{equation}
 f_i(\beta)=P_i(\beta)+aS_i[h_\beta(t_{M_i})-h_{\beta_0}(t_{M_i})],
 \label{eq:predictable-intercept-decomposition}
\end{equation}
with \(P_i\) measurable before \(S_i\). Removing only the current signed
pulse leaves an input bounded by \(U_*\), hence
\[
 |P_i(\beta)-P_i(\gamma)|\le U_*\|g_\beta-g_\gamma\|_1,
 \qquad |f_i(\beta)-f_i(\gamma)|\le d_*(\beta,\gamma).
\]
The full-input convolution also gives \(|f_i|\le2U_*C_*/\lambda_*\).
Use \(M=\max(1,2U_*C_*/\lambda_*)\) in the finite-n field estimates.

\begin{lemma}[Predictable-intercept information floor]
\label{lem:predictable-intercept-information}
Conditioning on every action selected before \(S_i\),
\[
 E[f_i(\beta)^2\mid\mathcal F_{i-1},r_i,u_i^0,E_i,M_i]
 =P_i(\beta)^2+a^2[h_\beta(t_{M_i})-h_{\beta_0}(t_{M_i})]^2.
\]
Consequently,
\begin{equation}
 E[f_i(\beta)^2\mid\mathcal F_{i-1}]\ge\rho a^2D(\beta,\beta_0).
 \label{eq:linear-time-information-floor}
\end{equation}
The same sign identity holds for directional derivatives on any fixed
finite coefficient cylinder.
\end{lemma}
\begin{proof}
The cross term has mean zero under the independent fresh sign. On an
exploration stage the remaining response square averages to \(a^2D\).
All other conditional-square terms are nonnegative. Differentiation of the
linear decomposition proves the cylinder statement.
\end{proof}

The exact initial-state contribution at stage \(i\) is
\(b_i^\beta(z)=\ell e^{\mathsf T_iA_\beta}z\), and is bounded by
\(r_i=C_*R e^{-\lambda_*\tau_gi}\). The finite-n correction lemma applies
with, explicitly,
\[
 R_1=\frac{C_*R}{e^{\lambda_*\tau_g}-1},\qquad
 R_2=\frac{C_*^2R^2}{e^{2\lambda_*\tau_g}-1}.
\]
Every fixed cylinder derivative is bounded by a polynomial in \(i\)
times this exponential, by the same-rate Duhamel calculation.

\begin{theorem}[Policy-uniform no-washout posterior exponent]
\label{thm:linear-time-adaptive-jacobi}
For each admissible policy sequence, with the coupling convention of
Section~\ref{sec:finite-sample-geometry}, almost surely for every closed
\(F\subset\mathfrak B\),
\begin{equation}
 \limsup_n n^{-1}\log\Pi_{n,\mathrm{lin}}^{\nu_n}(F)
 \le-\frac{\rho a^2}{2\sigma^2}\inf_{\beta\in F}D(\beta,\beta_0).
 \label{eq:linear-posterior-upper}
\end{equation}
For a coefficient cylinder the infimum is at least
\(\underline\kappa_J(\delta)\). In physical time the guaranteed upper
exponent is divided by at most \(\tau_+\). At every finite \(n\), the
field error is bounded by \eqref{eq:finite-n-field-bound} with the explicit
metric \(d_*\), and the exact-likelihood correction by
\eqref{eq:finite-n-transient-tail}, uniformly over all such policies.
\end{theorem}
\begin{proof}
All hypotheses of Proposition~\ref{prop:finite-n-field-concentration} have
just been checked with a deterministic compact metric common to every
history and policy. The conditional-square lower bound is
\eqref{eq:linear-time-information-floor}; the transient budget is the
summable global initial-state envelope above. Thus uniform likelihood
tracking \eqref{eq:adaptive-uniform-likelihood-tracking} holds, with
\(Q_n\ge\rho a^2D/(2\sigma^2)\).
To normalize, a \(d_*\)-ball of radius \(s\) about the truth satisfies
\(|f_i|\le s\) for all stages and policies. The finite-n denominator
argument of Theorem~\ref{thm:finite-n-posterior-separation}, or the uniform
Laplace bound with \(\varphi=0\), makes its exponent zero as
\(s\downarrow0\). Integrating the likelihood upper bound over \(F\)
proves the result. The effective inverse bound gives the cylinder claim.
Since \(\log\Pi(F)\le0\) and \(n/\mathsf T_n\ge1/\tau_+\), the
physical-time conclusion follows with the stated direction of inequality.
\end{proof}

The theorem is a guaranteed information floor. The actual adaptive
information \(Q_n\) also includes the baseline response. Rather than
identify it incorrectly with the floor, Theorem~\ref{thm:adaptive-laplace-tracking}
gives exact Laplace tracking and full LDPs along every uniformly convergent
information subsequence. A matching full-sequence LDP follows whenever
that predictable information converges uniformly. No convergence of the
information is assumed for the universal upper exponent.

For the washed experiment one may similarly draw a fresh exploration flag,
use \(\omega\) on exploration, and choose all other diagnostic cells
predictably before the fresh sign. With \(d=a\|h_\beta-h_\gamma\|_\infty\),
the same argument gives
\begin{equation}
 \limsup_n n^{-1}\log\Pi_n^{\nu_n}(F)
 \le-\frac{\rho a^2}{2\sigma^2}\inf_F D.
 \label{eq:adaptive-exploration-rate}
\end{equation}
This is the exploration-floor extension; the finite-n and growing-depth
versions are included in the next section.
